\documentclass{article}
\usepackage[T1]{fontenc}
\usepackage{lmodern}
\usepackage{graphicx}
\usepackage{amsmath,amssymb}
\usepackage[a4paper,margin=2cm]{geometry}
\usepackage[absolute,overlay]{textpos}
\setlength{\TPHorizModule}{1mm}
\setlength{\TPVertModule}{1mm}
\usepackage[indonesian]{babel}
\usepackage{hyperref}
\setlength{\emergencystretch}{2em}
% unit: Halaman 13

\begin{document}

% === Halaman 13 (python/native) ===

• Kunci untuk OTP harus seluruhnya acak sejati dan sepanjang pesan.

• Bagaimana jika kunci diambil dari teks yang panjang (misalnya tulisan 

di dalam novel, buku, berita, dan sebagainya)?

   - ini bukan lagi OTP  (sebab tulisan di buku/novel/berita bukan acak)

   - tidak menghasilkan perfect secrecy

   - dapat dipecahkan

• Kunci di dalam OTP hanya dipakai sekali dan tidak pernah digunakan 

kembali. Bagaimana jika kunci dipakai untuk kedua kalinya?

   - ia bukan lagi one-time pad, tetapi two-time pad

   - tidak aman 

13

\end{document}
